\subsection{Simulation}
\label{sec:simulation}
% 03_2_simulation.tex -- budget 0.5 p. C7; the planted spill-over control lives in 3.7 only.

We simulate sections from the generative model of \cref{sec:generative}: $6{,}000$ cells of eight spatially clustered types, a response linear in the neighbour composition, and spill-over planted at $\kappa = 0$ or $0.2$, three worlds each, each fitted with three seeds (\suppref{app:synthetic}, \cref{tab:synthetic}).

\paragraph{Recovery, and what the adversary does.}
The intrinsic state agrees with the planted type (NMI $0.80$ to $0.96$), as well as clusters of the counts or better in $17$ of $18$ fits; the response correlates with the planted one (first canonical correlation $0.83$ to $0.95$); the loadings span the planted programmes better than random subspaces in every fit; and $\vz$ carries almost none of the planted niche (within-type $R^2$ of composition at most $0.04$). Without the adversary, the response correlates at only $0.29$ to $0.91$ and the niche moves into $\vz$ ($R^2$ $0.11$ to $0.19$).

\paragraph{A wrong $\kappa$ costs little, so the fit cannot find the right one.}
On the worlds planted at $\kappa = 0.2$, fits at any assumed $\kappa$ from $0$ to $0.4$ recover the planted quantities about equally well (\cref{fig:synthetic-misspec}). Recovery does not hinge on guessing $\kappa$, and offers no handle on it: the non-identifiability of \cref{prop:kappa-bound}, seen in practice. The encoder's posterior mean is measurably off the exact one, more so at the final configuration than in development (\suppref{app:planted}). A positive control for the breakdown point, with spill-over planted at a known fraction, is in \cref{sec:robustness}.
